%% kapa: exam=1 round=29 year=2018 period=1 form=A subject=real_estate subject_ko=부동산학원론 seq=34 number=114
%% source: Q-Net 감정평가사 기출문제 (cst00309) artlSeq=5202808 file=extracted/1008/5202808/2018년 제29회 감정평가사 1차시험 기출문제/2018년 제29회 감정평가사 1차시험 기출문제/1교시 A형.pdf
%% conversion: kordoc 4.24.51 -> tools/gen1.py (auto); review=proofread-sonnet; flags=none
\begin{question}{114}
감정평가사 A는 표준지공시지가의 감정평가를 의뢰받고 현장조사를 통해 표준지에 대해 다음과 같이 확인하였다. 표준지조사평가보고서의 토지특성 기재방법으로 옳게 연결된 것은?
\begin{framedtext}
ㄱ. 토지이용상황: 주변의 토지이용상황이 ‘답’으로서 돈사ㆍ우사 등으로 이용되고 있는 토지 \par
ㄴ. 도로접면: 폭 12미터의 도로에 한면이 접하면서 자동차 통행이 가능한 폭 6미터의 도로에 다른 한면이 접하고 있는 토지 \par
\end{framedtext}
\begin{choices}
  \item ㄱ: 목장용지, ㄴ: 중로각지
  \item ㄱ: 목장용지, ㄴ: 소로각지
  \item ㄱ: 답기타, ㄴ: 중로각지
  \item ㄱ: 답기타, ㄴ: 소로각지
  \item ㄱ: 답축사, ㄴ: 중로각지
\end{choices}
\end{question}
\answer{5}
